\section{GPU-resident meshing: design constraints}
\label{sec:gpudesign}

\subsection{The contract}

The GPU mesher must satisfy a strict contract:
\begin{enumerate}[label=\textbf{G\arabic*.},leftmargin=2.2em]
\item \textbf{No host vertex stream.}  The host may upload node tables,
  scan-cell indices, and small counters.  It must not upload or download a
  Cartesian vertex.
\item \textbf{Renderer layout.}  The output vertex must be the renderer's
  64-byte \code{Vertex}; the output index must be a chunk-local slot index
  that the renderer's draw command can rebase with \code{vertexOffset}.
\item \textbf{Bit-exact parity with the CPU reference.}  The same chunk
  must produce the same triangle set and the same QEF positions as the CPU
  walker, up to the documented floating-point ordering.
\item \textbf{Deterministic slots.}  The host must assign vertex slots in a
  stable order so a rebuild produces the same pool layout and the
  ray-tracing BLAS does not see gratuitous changes.
\item \textbf{Bounded work and memory.}  A pathological chunk must fail
  explicitly (overflow/drop counter) rather than corrupt the pool or draw a
  partial surface.
\end{enumerate}

\subsection{Why one geometry pass?}

The walk output is not known before the scan: any active edge can emit any
number of sub-segments, and each sub-segment expands to one or two
triangles.  A two-stage design solved this by writing one raw node key per
triangle corner, reading the count back, reserving the exact index span,
and remapping the keys in a second dispatch.  That intermediate stream is
unnecessary once the host can bound the output from the topology it already
has.

While the scan-cell list is built, the host counts the chunk's
\emph{cross edges}: the \code{(cell, cube-edge)} pairs whose two stored
corner samples have opposite signs.  The shader scans exactly those edges,
so the count matches its work list.  The index capacity is then a topology
bound, not a calibrated multiplier:
\begin{equation}
  I_{\text{cap}} =
  \sum_{\text{scan cell }c}
    6 \cdot \text{segments}(c) \cdot \text{crossEdges}(c),
  \label{eq:indexcap}
\end{equation}
where a segment expands to at most two triangles of three indices and
\begin{equation}
  \text{segments}(c) =
  \operatorname{clamp}\!\left(
    \frac{\text{cellSize}(c)}{\text{smallestFlattenedCellSize}},
    1, 64\right).
  \label{eq:segments}
\end{equation}
Every break coordinate the collector can insert is a boundary of a cell
returned by \code{findCellAt}; all such cubes are aligned to the smallest
cube in the flattened table, so one edge of size $S$ can contain at most
$S/m$ dyadic intervals.  The clamp is the shader's own break-list capacity.

The vertex capacity is exactly the dense surface-node slot address space
\code{activeVertexCount}.  The host assigns one slot per surface node in
table order, and every index is one of those slots.  The shader writes only
the slots a triangle references, so the published extent can be smaller
than the reserved capacity, but the address space itself is an exact upper
bound.

\begin{figure}[htbp]
\centering
\resizebox{\linewidth}{!}{%
\begin{tikzpicture}[node distance=6mm and 9mm,font=\small]
  \node[io] (in) {flattened node table\\+ compacted scan cells};
  \node[cpu,right=of in] (cap) {CPU cross-edge\\capacity plan};
  \node[io,right=of cap] (host) {reserve vertex/index\\spans in packed pools};
  \node[gpu,right=of host] (pass) {geometry pass\\\code{main}};
  \node[io,right=of pass] (out) {packed pools\\index + 64B vertex};
  \draw[arrow] (in)--(cap);
  \draw[arrow] (cap)--(host);
  \draw[arrow] (host)--(pass);
  \draw[arrow] (pass)--(out);
  \draw[darrow] (pass.south) -- ++(0,-6mm) -| node[pos=0.25,below,note]{overflow: exact counts, release + re-reserve} (host.south);
\end{tikzpicture}}
\caption{The single geometry dispatch.  The CPU sizes both spans from the
cross-edge topology before submission; the shader appends indices directly
and writes each referenced vertex payload once.  A defensive overflow is
reported through counters and resolved by growing only this chunk's span.}
\label{fig:geometrypass}
\end{figure}

\subsection{Overflow without out-of-bounds writes}

The shader still reserves an index base with \code{atomicAdd} and writes the
triangle only when
\begin{equation}
  \text{base}+3 \le I_{\text{cap}}
  \quad\land\quad
  \text{extent} \le V_{\text{cap}}.
\end{equation}
On failure it increments an overflow counter and drops the triangle, but it
has already advanced the index counter by three and updated the maximum
referenced slot.  The host therefore reads back the exact index requirement
and vertex extent without any write outside the reserved span.  It releases
that chunk's span and asks the provider for the grown capacities; the retry
is limited to this chunk and normally the first grown attempt succeeds.  A
fixed-array overflow or a corner without a host-assigned slot is not
retryable: that indicates a flattening disagreement, so the chunk is
failed.

\subsection{Determinism}

The walker uses atomics for the index append, the overflow counters, and
the per-node vertex election.  Triangle order at the end of the index
stream is therefore non-deterministic when two scan cells emit
concurrently.  Does that matter?  The rendered triangle \emph{set} is
deterministic, and each corner carries its host-assigned slot.  The
renderer does not care about triangle order.  The vertex slots are assigned
by the host in node-table order, not by emission order, so a rebuild
produces the same vertex buffer for the same node table.  Ray tracing and
shadow maps are order-independent.  The only observable order is the
index stream, which is internal to the chunk.

\subsection{Precision strategy}

The shader is a direct translation of the CPU arithmetic:
\begin{itemize}
\item corner order, edge order, and QEF accumulation order are identical;
\item the trilinear lerp uses the same guarded branch;
\item the QEF uses the same Cramer solve and the same $10^{-2}$ regularizer;
\item the gradient and QEF use \code{precise} to prevent FMA contraction;
\item the sample nudge uses \code{nextafter} semantics, including the
  subnormal guard at zero.
\end{itemize}
The goal is not merely a visually similar mesh; it is a mesh that the
offline exporter can compare triangle-by-triangle.  Where the GPU cannot
guarantee the same order (atomic append), the triangle set is still
identical because the walker's per-cell output is independent.
